\documentclass[a4paper]{article}
\usepackage[a4paper,margin=0mm]{geometry}
\usepackage{fontspec}
\usepackage{unicode-math}
\usepackage{tikz}
\usepackage[table]{xcolor}
\usepackage[absolute]{textpos}
\usepackage{array}
\usepackage{multirow}

\setmainfont{Libertinus Serif}
\setsansfont{Libertinus Sans}
\setmathfont{Libertinus Math}

\definecolor{accent}{HTML}{0F6A52}
\definecolor{badge}{HTML}{A9CFC2}
\definecolor{ink}{HTML}{2A2E30}

\setlength{\TPHorizModule}{1mm}
\setlength{\TPVertModule}{1mm}
\setlength{\parindent}{0pt}
\setlength{\parskip}{0pt}
\emergencystretch=3em
\abovedisplayskip=4pt plus 2pt minus 2pt \belowdisplayskip=4pt plus 2pt minus 2pt \abovedisplayshortskip=2pt \belowdisplayshortskip=2pt
\pagestyle{empty}

\newcommand{\LCOL}{16mm}
\newcommand{\RCOL}{111mm}
\newcommand{\COLW}{83mm}

\newcommand{\topicheader}[2]{%
\begin{textblock*}{194mm}(16mm,10mm)
\begin{tikzpicture}
\node[fill=accent,text=white,font=\sffamily\bfseries\fontsize{13}{13}\selectfont,
      minimum width=9mm,minimum height=9mm,inner sep=0pt] at (4.5mm,-4.5mm) {#1};
\node[font=\sffamily\bfseries\fontsize{15}{16}\selectfont,text=ink,anchor=west] at (12mm,-4.5mm) {#2};
\end{tikzpicture}
\end{textblock*}
\begin{textblock*}{194mm}(16mm,22.5mm)
\begin{tikzpicture}
\fill[accent] (0,0) rectangle (194mm,0.5mm);
\end{tikzpicture}
\end{textblock*}
}

\newcommand{\pagefooter}[1]{%
\begin{textblock*}{194mm}(8mm,286.41mm)
  \begin{tikzpicture}
    \fill[accent] (0,0) rectangle (194mm,0.6mm);
  \end{tikzpicture}
\end{textblock*}
\begin{textblock*}{63.59mm}(8mm,288.27mm)
  {\sffamily\bfseries\color{accent}\fontsize{13}{16}\selectfont SAT 1200+}
\end{textblock*}
\begin{textblock*}{10mm}(188mm,288.27mm)
  \centering{\sffamily\color{accent}\fontsize{14}{16}\selectfont #1}
\end{textblock*}
}

\newcommand{\colrule}[2]{%
\begin{textblock*}{1mm}(105mm,#1mm)
\begin{tikzpicture}
\draw[black!65,line width=0.5pt] (0,0) -- (0,#2mm);
\end{tikzpicture}
\end{textblock*}
}

\newcommand{\setkicker}[2]{%
{\sffamily\bfseries\fontsize{12}{14}\selectfont #1.}\ {\itshape\fontsize{11}{14}\selectfont\color{black!75} #2}\par\vspace{2mm}%
}

\setlength{\fboxsep}{1.4pt}
\setlength{\fboxrule}{0.5pt}
\newcommand{\pnum}[1]{\fcolorbox{ink}{white}{\sffamily\bfseries\fontsize{8}{8}\selectfont #1}~}

% ---- SAT-style multiple-choice block (2x2 grid of short choices) ----
\newcommand{\satq}[6]{%
\noindent\pnum{#1} #2\par\nobreak\vspace{2mm}%
\noindent\hspace*{7mm}\begin{tabular}[t]{@{}p{32mm}p{32mm}@{}}
(A) #3 & (B) #4 \\[2mm]
(C) #5 & (D) #6 \\
\end{tabular}\par\vspace{5.5mm}%
}

% ---- stacked choices for long ones ----
\newcommand{\satqw}[6]{%
\noindent\pnum{#1} #2\par\nobreak\vspace{2mm}%
\noindent\hspace*{7mm}(A) #3\par\vspace{1.6mm}%
\noindent\hspace*{7mm}(B) #4\par\vspace{1.6mm}%
\noindent\hspace*{7mm}(C) #5\par\vspace{1.6mm}%
\noindent\hspace*{7mm}(D) #6\par\vspace{5.5mm}%
}

% ---- student-produced response (grid-in) ----
\newcommand{\satg}[2]{%
\noindent\pnum{#1} {\sffamily\bfseries\fontsize{7}{8}\selectfont\color{accent}GRID-IN}\ #2\par\nobreak\vspace{2.5mm}%
\noindent\hspace*{7mm}{\sffamily\fontsize{8.5}{10}\selectfont Answer:}\ \rule[-1.2mm]{28mm}{0.4pt}\par\vspace{5.5mm}%
}

% ---- worked-solution entry ----
\newcommand{\solitem}[4]{%
\noindent\pnum{#1} {\bfseries #2}\ —\ #3\ {\itshape\color{accent}[#4]}\par\vspace{2.1mm}%
}

\arrayrulecolor{accent}
\setlength{\arrayrulewidth}{0.7pt}
\renewcommand{\arraystretch}{1.3}

\usepackage{multicol}
\usepackage{fancyhdr}
\usepackage{needspace}
\newcommand{\pgoff}{0}
\renewcommand{\solitem}[4]{\noindent\parbox[t]{\linewidth}{\raggedright\pnum{#1} {\bfseries #2}\ ---\ #3\ {\itshape\color{accent}[#4]}}\par\vspace{3.6mm}}
\newcommand{\solkicker}[2]{\Needspace{38mm}\vspace{1.5mm}{\sffamily\bfseries\fontsize{12.5}{15}\selectfont #1.}\ {\itshape\fontsize{11}{14}\selectfont\color{black!75} #2}\par\vspace{3mm}}
\renewcommand{\satqw}[6]{%
\noindent\pnum{#1} #2\par\nobreak\vspace{2mm}%
\noindent\hangindent=13mm\hangafter=1\hspace*{7mm}(A)~#3\par\vspace{1.6mm}%
\noindent\hangindent=13mm\hangafter=1\hspace*{7mm}(B)~#4\par\vspace{1.6mm}%
\noindent\hangindent=13mm\hangafter=1\hspace*{7mm}(C)~#5\par\vspace{1.6mm}%
\noindent\hangindent=13mm\hangafter=1\hspace*{7mm}(D)~#6\par\vspace{5.5mm}%
}
\newcommand{\solkitem}[5]{\noindent\parbox[t]{\linewidth}{\vspace{1.5mm}{\sffamily\bfseries\fontsize{12.5}{15}\selectfont #1.}\ {\itshape\fontsize{11}{14}\selectfont\color{black!75} Worked solutions}\par\vspace{3mm}\raggedright\pnum{#2} {\bfseries #3}\ ---\ #4\ {\itshape\color{accent}[#5]}}\par\vspace{3.6mm}}
\fancypagestyle{solstyle}{\fancyhf{}\renewcommand{\headrulewidth}{0pt}\renewcommand{\footrulewidth}{0pt}%
 \fancyhead[L]{\topicheader{2}{Hard Linear Systems}}%
 \fancyfoot[L]{\pagefooter{\the\numexpr\value{page}+\pgoff\relax}}}

\begin{document}
\topicheader{2}{Hard Linear Systems}
\pagefooter{1}
\colrule{32}{238}
\mbox{}%
\begin{textblock*}{\COLW}(\LCOL,32mm)
{\fontsize{10}{12.5}\selectfont
\setkicker{Set 1}{Bridge: the core moves, one problem each.}
\satq{1}{In the system $kx+12y=9$ and $5x+4y=7$, $k$ is a constant. For which value of $k$ does the system have no solution?}{$\dfrac53$}{$3$}{$15$}{$20$}
\satg{2}{In the system $2x+7y=b$ and $ax+21y=30$, $a$ and $b$ are constants. The system has infinitely many solutions. What is the value of $a+b$?}
\satq{3}{One equation of a system is $4(x-2y)=3y+10$. The system has no solution. Which could be the second equation?}{$-8x+22y=-5$}{$8x-22y=20$}{$11x-4y=10$}{$4x+11y=10$}
\satq{4}{At a bakery, 5 muffins and 4 bagels cost \$21.00, and 3 muffins and 4 bagels cost \$15.00. What is the cost of one bagel?}{$\$6.00$}{$\$4.00$}{$\$3.00$}{$\$1.50$}
\satqw{5}{A teacher spends exactly \$600 on $x$ notebooks at \$3 each and $y$ binders at \$5 each, so $3x+5y=600$. What does the point $(0,120)$ on the graph of this equation represent?}{120 notebooks and no binders cost a total of \$600.}{No notebooks and 120 binders cost a total of \$600.}{The teacher can buy 120 items of either kind for \$600.}{The teacher buys 120 binders for each notebook.}
\satq{6}{Line $\ell$ passes through $(2,5)$ and $(6,13)$. Line $m$ is parallel to $\ell$ and passes through $(1,-4)$. What is the $y$-intercept of line $m$?}{$-6$}{$-4$}{$2$}{$6$}
}
\end{textblock*}
\begin{textblock*}{\COLW}(\RCOL,32mm)
{\fontsize{10}{12.5}\selectfont
\setkicker{Set 1 (cont.)}{Bridge}
\satq{7}{Line $q$ is perpendicular to the line $2x+3y=18$ and passes through $(10,6)$. What is the $x$-intercept of $q$?}{$1$}{$6$}{$14$}{$19$}
\satq{8}{What is the sum of the solutions of $|4x-6|=14$?}{$7$}{$5$}{$3$}{$-2$}
\satq{9}{For a constant $k$, the equation $3(2x-k)+x=7x+12$ has infinitely many solutions. What is the value of $k$?}{$-12$}{$-4$}{$4$}{$12$}
\satg{10}{What is the greatest integer $x$ that satisfies $3(2x-5)<4x+9$?}
\satq{11}{Which ordered pair $(x,y)$ satisfies both $y>2x-3$ and $x+y<6$?}{$(0,6)$}{$(4,1)$}{$(6,3)$}{$(2,3)$}
\satg{12}{A linear function $f$ has $f(-2)=19$ and $f(6)=-13$. For what value of $x$ does $f(x)=0$?}
\vspace{3mm}
\setkicker{Set 2}{Parameters: when are there no, one or infinitely many solutions?}
\satq{1}{In the system $kx+9y=6$ and $4x+ky=-4$, $k$ is a constant. For which value of $k$ does the system have no solution?}{$-6$}{$\dfrac32$}{$6$}{$36$}
\satg{2}{In the system $\dfrac13x+\dfrac45y=2$ and $px+qy=9$, $p$ and $q$ are constants and the system has infinitely many solutions. What is the value of $p+q$?}
}
\end{textblock*}
\clearpage
\topicheader{2}{Hard Linear Systems}
\pagefooter{2}
\colrule{32}{238}
\mbox{}%
\begin{textblock*}{\COLW}(\LCOL,32mm)
{\fontsize{10}{12.5}\selectfont
\setkicker{Set 2 (cont.)}{Parameters}
\satq{3}{One equation of a system is $5x-2(y+3)=2y+14$. The system has no solution. Which could be the second equation?}{$-10x+8y=-5$}{$10x-8y=40$}{$4x+5y=20$}{$5x+4y=20$}
\satq{4}{Which equation, together with $\dfrac25x-\dfrac13y=2$, forms a system with infinitely many solutions?}{$5x-6y=30$}{$6x-5y=15$}{$6x+5y=30$}{$6x-5y=30$}
\satq{5}{In the system $ax+by=20$ and $5x+2by=18$, $a$ and $b$ are constants. The graphs of the equations in the $xy$-plane intersect at the point $(2,y)$. What is the value of $a$?}{$10$}{$8$}{$6$}{$5.5$}
\satq{6}{If $5a+2b=34$ and $3a-4b=10$, what is the value of $a+3b$?}{$44$}{$22$}{$12$}{$6$}
\satq{7}{$6x+4y=50$ and $4x+6y=40$. What is the mean (average) of $x$ and $y$?}{$4.5$}{$9$}{$45$}{$90$}
\satq{8}{For how many real values of $k$ does the system $kx+4y=7$ and $x+ky=3$ have no solution?}{$4$}{$2$}{$1$}{$0$}
\satg{9}{The graphs of $2x+3y=k$ and $x-y=2$ in the $xy$-plane intersect at a point on the line $y=2x$. What is the value of $k$?}
}
\end{textblock*}
\begin{textblock*}{\COLW}(\RCOL,32mm)
{\fontsize{10}{12.5}\selectfont
\setkicker{Set 2 (cont.)}{Parameters}
\satq{10}{In the system $ax-6y=18$ and $2x+3y=b$, $a$ and $b$ are constants. The system has infinitely many solutions. What is the value of $a-b$?}{$-13$}{$-5$}{$5$}{$36$}
\satg{11}{Line $\ell$ has equation $y=mx+b$. The system formed by $\ell$ and $3x-2y=5$ has no solution, and $\ell$ passes through $(4,7)$. What is the value of $b$?}
\satqw{12}{Which of the following systems of equations has infinitely many solutions?}{$2x+6y=10$ and $x+3y=10$}{$5x+y=4$ and $10x+2y=6$}{$x-4y=2$ and $4x-y=8$}{$3x-2y=8$ and $-6x+4y=-16$}
\satg{13}{In the system $y=3x+k$ and $2y=6x+k+9$, $k$ is a constant. For one value of $k$ the system has infinitely many solutions. What is that value of $k$?}
\satq{14}{In the system $y=2x+a$ and $y=-x+b$, $a$ and $b$ are constants. The solution $(x,y)$ of the system has $x>0$. Which of the following must be true?}{$a<b$}{$a>b$}{$a+b>0$}{$ab>0$}
\vspace{3mm}
\setkicker{Set 3}{Context: build the equations, then interpret the numbers.}
\satq{1}{A chemist mixes a $20\%$ acid solution with a $50\%$ acid solution to make 90 liters of a $30\%$ acid solution. How many liters of the $50\%$ solution are used?}{$27$}{$30$}{$45$}{$60$}
}
\end{textblock*}
\clearpage
\topicheader{2}{Hard Linear Systems}
\pagefooter{3}
\colrule{32}{238}
\mbox{}%
\begin{textblock*}{\COLW}(\LCOL,32mm)
{\fontsize{10}{12.5}\selectfont
\setkicker{Set 3 (cont.)}{Context}
\satg{2}{A boat travels 36 kilometers upstream in 3 hours. The return trip downstream along the same route takes 2 hours. The boat's speed in still water and the speed of the current are both constant. What is the speed of the current, in kilometers per hour?}
\satqw{3}{A van carries $x$ boxes of 18 kilograms each and $y$ crates of 30 kilograms each, with a total load of exactly 450 kilograms: $18x+30y=450$. If the number of boxes increases by 5 and the total load stays 450 kilograms, how must the number of crates change?}{It increases by $3$.}{It decreases by $\frac35$.}{It decreases by $3$.}{It decreases by $\frac{25}{3}$.}
\satq{4}{The table gives values of the linear function $f$. $$\begin{array}{c|ccc}x&2&5&9\\\hline f(x)&9&18&30\end{array}$$ The graph of $g$ is the graph of $f$ shifted 4 units to the right. What is the $y$-intercept of the graph of $g$?}{$15$}{$7$}{$-1$}{$-9$}
\satq{5}{A courier charges a flat fee for any parcel up to 2 kilograms, plus $c$ dollars for each kilogram above 2. A 4-kilogram parcel costs \$26 and a 7-kilogram parcel costs \$44. What is the cost of a 10-kilogram parcel?}{$\$56$}{$\$60$}{$\$62$}{$\$74$}
\satq{6}{On Monday a cafe sold 30 small and 20 large coffees for \$190. On Tuesday it sold 20 small and 30 large coffees for \$210. At these prices, how much money, in dollars, would 10 small and 10 large coffees cost?}{$8$}{$80$}{$200$}{$400$}
}
\end{textblock*}
\begin{textblock*}{\COLW}(\RCOL,32mm)
{\fontsize{10}{12.5}\selectfont
\setkicker{Set 3 (cont.)}{Context}
\satqw{7}{A tank is filled by two pipes. If pipe A runs for $x$ minutes and pipe B runs for $y$ minutes, the tank is exactly full when $\dfrac{x}{12}+\dfrac{y}{20}=1$. What does the number 12 represent?}{The number of minutes pipe A needs to fill the tank by itself.}{The number of minutes both pipes together need to fill the tank.}{The fraction of the tank that pipe A fills in one minute.}{The number of minutes pipe B needs to fill the tank by itself.}
\satg{8}{A cyclist rides for a total of 3.5 hours, going 18 kilometers per hour on flat roads and 12 kilometers per hour uphill, and covers 54 kilometers. For how many minutes does the cyclist ride uphill?}
\satg{9}{A linear model predicts the number of subscribers, in thousands, of a streaming service $t$ months after launch. The model gives 52 thousand at 2 months and 88 thousand at 5 months. According to the model, what is the least whole number of months after launch at which the subscribers exceed 120 thousand?}
\satq{10}{For a linear function $f$ and a constant $a$, $f(a+3)-f(a)=12$. What is the value of $f(a+10)-f(a)$?}{$4$}{$22$}{$40$}{$120$}
\satg{11}{On a test, each correct answer earns 4 points, each incorrect answer loses 1 point, and each unanswered question earns 0 points. Dana answered exactly 30 questions and scored 90 points. How many of her answers were incorrect?}
\satq{12}{Maya is 4 years older than twice Jon's age. In 6 years, the sum of their ages will be 52. How old is Maya now?}{$12$}{$28$}{$34$}{$36$}
}
\end{textblock*}
\clearpage
\topicheader{2}{Hard Linear Systems}
\pagefooter{4}
\colrule{32}{238}
\mbox{}%
\begin{textblock*}{\COLW}(\LCOL,32mm)
{\fontsize{10}{12.5}\selectfont
\setkicker{Set 3 (cont.)}{Context}
\satq{13}{Plan A costs \$28 per month plus \$0.12 per gigabyte (GB) used. Plan B costs \$18 per month plus \$0.20 per GB used. For how many GB of monthly use do the two plans cost the same?}{$12.5$}{$50$}{$125$}{$575$}
\satqw{14}{The linear functions $f$ and $g$ are defined by $f(x)=2x+3$ and $g(x)=-x+5$. For how many values of $x$ is $f(g(x))=g(f(x))$?}{Infinitely many}{$2$}{$1$}{$0$}
\vspace{3mm}
\setkicker{Set 4}{Inequalities, absolute value and identities.}
\satq{1}{What is the sum of all solutions of $|x-4|=2x-5$?}{$4$}{$3$}{$1$}{$0$}
\satg{2}{How many integers $x$ satisfy $|2x-7|\le9$?}
\satq{3}{A bakery bakes $x$ loaves and $y$ rolls in a batch. The oven holds at most 50 items, so $x+y\le50$. Loaves sell for \$6 and rolls for \$4, and the batch must bring in at least \$220, so $6x+4y\ge220$. Which ordered pair $(x,y)$ satisfies both inequalities?}{$(30,15)$}{$(10,35)$}{$(25,15)$}{$(35,20)$}
\satg{4}{A course grade is $70\%$ of the test average plus $30\%$ of the final exam score. Priya's test average is 74. What is the least final exam score that gives her a course grade of at least 80?}
}
\end{textblock*}
\begin{textblock*}{\COLW}(\RCOL,32mm)
{\fontsize{10}{12.5}\selectfont
\setkicker{Set 4 (cont.)}{Inequalities}
\satq{5}{The solution set of $ax-12<5x+6$ is $x>-3$, where $a$ is a constant. What is the value of $a$?}{$11$}{$7$}{$-1$}{$-11$}
\satg{6}{For constants $k$ and $m$, the equation $\dfrac{kx-9}{3}=5x+m$ has infinitely many solutions. What is the value of $k+m$?}
\satqw{7}{How many solutions does the following equation have? $$5(x-2)-3(x+4)=2x-20$$}{Infinitely many}{Exactly two}{Exactly one}{No solution}
\satq{8}{In the $xy$-plane, the line $kx+3y=12$ is perpendicular to the line $4x-6y=7$, where $k$ is a constant. What is the value of $k$?}{$-4.5$}{$-2$}{$2$}{$4.5$}
\satq{9}{The lines $y=3x+2$ and $y=-x+10$ and the $y$-axis form a triangle. What is the area of the triangle?}{$8$}{$12$}{$16$}{$32$}
\satq{10}{A truck can carry at most 2,400 pounds. A shipment contains small crates of 30 pounds and large crates of 75 pounds, with at least twice as many small crates as large crates. What is the greatest number of large crates the truck can carry?}{$17$}{$18$}{$24$}{$32$}
\satg{11}{The point $(k,6)$ satisfies both $y<3x-3$ and $x+y\le13$. How many integer values of $k$ are possible?}
}
\end{textblock*}
\clearpage
\topicheader{2}{Hard Linear Systems}
\pagefooter{5}
\colrule{32}{238}
\mbox{}%
\begin{textblock*}{\COLW}(\LCOL,32mm)
{\fontsize{10}{12.5}\selectfont
\setkicker{Set 4 (cont.)}{Inequalities}
\satq{12}{For real numbers $x$ and $y$, $|x-6|\le4$ and $|y+2|\le3$. What is the greatest possible value of $x-y$?}{$9$}{$11$}{$13$}{$15$}
\satqw{13}{A part's diameter $d$, in millimeters, is acceptable if $14.2\le d\le14.8$. Which inequality describes exactly the acceptable diameters?}{$|d-14.5|\le0.6$}{$|d-14.5|\le0.3$}{$|d-14.8|\le0.3$}{$|d-0.3|\le14.5$}
\satqw{14}{How many solutions does the equation $|x+2|+|x-4|=6$ have?}{No solution}{Exactly one}{Exactly two}{Infinitely many}
\vspace{3mm}
\setkicker{Set 5}{Summit: the hardest mixed problems.}
\satq{1}{In the system $kx+2y=10$ and $3x+(k-1)y=10$, $k$ is a constant. For which value of $k$ does the system have no solution?}{$-3$}{$-2$}{$2$}{$3$}
\satg{2}{Line $\ell$ passes through $(2k,5)$ and $(6k,17)$, where $k$ is a positive constant, and $\ell$ is parallel to the line $3x-2y=8$. What is the $y$-intercept of $\ell$?}
\satq{3}{For which value of the constant $c$ does the equation $2(cx-3)=4x-c$ have no solution?}{$6$}{$4$}{$2$}{$-2$}
}
\end{textblock*}
\begin{textblock*}{\COLW}(\RCOL,32mm)
{\fontsize{10}{12.5}\selectfont
\setkicker{Set 5 (cont.)}{Summit}
\satg{4}{The numbers $x$ and $y$ satisfy $2x+y=3k$ and $x-2y=k+5$, where $k$ is a constant. If $x+y=7$, what is the value of $k$?}
\satq{5}{How many liters of pure water must be added to 40 liters of a $35\%$ salt solution to make a $28\%$ salt solution?}{$7$}{$8$}{$10$}{$14$}
\satg{6}{How many ordered pairs $(x,y)$ of positive integers satisfy $2x+3y\le20$?}
\satq{7}{The function $f(x)=ax+b$ has $a>0$ and satisfies $f(f(x))=9x+8$. What is the value of $f(1)$?}{$11$}{$5$}{$3$}{$-7$}
\satq{8}{How many solutions does the equation $\big||x-3|-2\big|=1$ have?}{$4$}{$3$}{$2$}{$1$}
\satq{9}{A school orders $x$ adult tickets at \$5 each and $y$ student tickets at \$3 each. It needs at least 12 tickets in all, at least 3 adult tickets, and at most 8 student tickets. What is the least possible total cost, in dollars?}{$36$}{$39$}{$44$}{$60$}
\satg{10}{Car A leaves a rest stop at 1:00 p.m. and travels at 48 miles per hour. Car B leaves the same rest stop at 1:45 p.m. on the same road and travels at 60 miles per hour. How many miles from the rest stop is Car B when it catches up to Car A?}
}
\end{textblock*}
\clearpage
\topicheader{2}{Hard Linear Systems}
\pagefooter{6}
\colrule{32}{238}
\mbox{}%
\begin{textblock*}{\COLW}(\LCOL,32mm)
{\fontsize{10}{12.5}\selectfont
\setkicker{Set 5 (cont.)}{Summit}
\satqw{11}{Plan A costs $C=40+12m$ dollars after $m$ months, and Plan B costs $C=100+7m$ dollars. The system has the solution $(12,184)$. Which statement is true?}{Plan A costs more than Plan B for every $m$.}{Plan B costs less than Plan A when $m<12$.}{Plan A costs less than Plan B when $m>12$.}{Plan A costs more than Plan B when $m>12$.}
\satq{12}{The graphs of $y=ax+3$ and $y=bx-5$, where $a$ and $b$ are constants, intersect at the point whose $x$-coordinate is 4. What is the value of $a-b$?}{$-8$}{$-2$}{$2$}{$8$}
\satq{13}{Pipe A fills an empty tank in 6 hours and pipe B fills it in 9 hours. Drain C empties a full tank in 18 hours. With all three open, how many hours does it take to fill the empty tank?}{$3.6$}{$4.5$}{$7.5$}{$18$}
\satq{14}{Line $m$ is perpendicular to the line $4x-3y=12$ and has the same $y$-intercept as the line $2x+5y=20$. Which is an equation of line $m$?}{$3x+4y=16$}{$3x-4y=-16$}{$4x+3y=12$}{$3x+4y=-16$}
\satq{15}{The equation $|3x-a|=12$ has two solutions whose sum is 10, where $a$ is a constant. What is the value of $a$?}{$5$}{$10$}{$15$}{$30$}
\satg{16}{How many integers $x$ satisfy $-7<\dfrac{5-2x}{3}\le5$?}
\satq{17}{Class A has 24 students with a mean score of 78. Class B has 36 students. The mean score of all 60 students is 84. What is the mean score of Class B?}{$84$}{$88$}{$90$}{$93$}
}
\end{textblock*}
\begin{textblock*}{\COLW}(\RCOL,32mm)
{\fontsize{10}{12.5}\selectfont
\setkicker{Set 5 (cont.)}{Summit}
\satq{18}{For which values of $m$ do the graphs of $y=mx+3$ and $y=-x+9$ intersect in Quadrant I?}{$m>-1$}{$m>0$}{$m<-1$}{$m>-\frac13$}
\satq{19}{If $\dfrac{2x+y}{x-3y}=\dfrac52$ and $y\ne0$, what is the value of $\dfrac xy$?}{$17$}{$\dfrac1{17}$}{$-13$}{$-17$}
\satg{20}{The table gives two points on a line in the $xy$-plane. $$\begin{array}{c|cc}x&p&p+4\\\hline y&2p+1&2p+17\end{array}$$ The $y$-intercept of the line is $(0,-7)$. What is the value of $p$?}
}
\end{textblock*}
\clearpage
\topicheader{2}{Hard Linear Systems}
\pagefooter{7}
\mbox{}%
\begin{textblock*}{178mm}(16mm,30mm)
{\bfseries\itshape\color{accent}\fontsize{19}{21}\selectfont\centering Answer Key\par}
\end{textblock*}
\begin{textblock*}{178mm}(16mm,41mm)
{\fontsize{9.5}{12}\selectfont\color{black!70}\centering Multiple choice: the letter. Grid-in: the number to enter. Worked solutions follow on the next pages.\par}
\end{textblock*}
\begin{textblock*}{178mm}(16mm,52mm)
{\fontsize{10}{13}\selectfont
\centering
\setlength{\tabcolsep}{2pt}\renewcommand{\arraystretch}{1.3}
{\sffamily\bfseries\color{accent}\fontsize{11}{13}\selectfont Set 1 \ \textmd{\itshape\small Bridge: the core moves}}\\[1.3mm]
\begin{tabular}{|>{\centering\arraybackslash}p{16mm}|>{\centering\arraybackslash}p{16mm}|>{\centering\arraybackslash}p{16mm}|>{\centering\arraybackslash}p{16mm}|>{\centering\arraybackslash}p{16mm}|>{\centering\arraybackslash}p{16mm}|>{\centering\arraybackslash}p{16mm}|>{\centering\arraybackslash}p{16mm}|>{\centering\arraybackslash}p{16mm}|>{\centering\arraybackslash}p{16mm}|}
\hline
\cellcolor{badge!55}\textbf{1} & \cellcolor{badge!55}\textbf{2} & \cellcolor{badge!55}\textbf{3} & \cellcolor{badge!55}\textbf{4} & \cellcolor{badge!55}\textbf{5} & \cellcolor{badge!55}\textbf{6} & \cellcolor{badge!55}\textbf{7} & \cellcolor{badge!55}\textbf{8} & \cellcolor{badge!55}\textbf{9} & \cellcolor{badge!55}\textbf{10}\\
\hline
\textbf{C} & 16 & \textbf{A} & \textbf{D} & \textbf{B} & \textbf{A} & \textbf{B} & \textbf{C} & \textbf{B} & 11\\
\hline
\end{tabular}\\[1.5mm]
\begin{tabular}{|>{\centering\arraybackslash}p{16mm}|>{\centering\arraybackslash}p{16mm}|}
\hline
\cellcolor{badge!55}\textbf{11} & \cellcolor{badge!55}\textbf{12}\\
\hline
\textbf{D} & 11/4\\
\hline
\end{tabular}\\[1.5mm]
\vspace{3.5mm}
{\sffamily\bfseries\color{accent}\fontsize{11}{13}\selectfont Set 2 \ \textmd{\itshape\small Systems with parameters: none, one or infinitely many}}\\[1.3mm]
\begin{tabular}{|>{\centering\arraybackslash}p{16mm}|>{\centering\arraybackslash}p{16mm}|>{\centering\arraybackslash}p{16mm}|>{\centering\arraybackslash}p{16mm}|>{\centering\arraybackslash}p{16mm}|>{\centering\arraybackslash}p{16mm}|>{\centering\arraybackslash}p{16mm}|>{\centering\arraybackslash}p{16mm}|>{\centering\arraybackslash}p{16mm}|>{\centering\arraybackslash}p{16mm}|}
\hline
\cellcolor{badge!55}\textbf{1} & \cellcolor{badge!55}\textbf{2} & \cellcolor{badge!55}\textbf{3} & \cellcolor{badge!55}\textbf{4} & \cellcolor{badge!55}\textbf{5} & \cellcolor{badge!55}\textbf{6} & \cellcolor{badge!55}\textbf{7} & \cellcolor{badge!55}\textbf{8} & \cellcolor{badge!55}\textbf{9} & \cellcolor{badge!55}\textbf{10}\\
\hline
\textbf{C} & 5.1 & \textbf{A} & \textbf{D} & \textbf{B} & \textbf{C} & \textbf{A} & \textbf{B} & −16 & \textbf{C}\\
\hline
\end{tabular}\\[1.5mm]
\begin{tabular}{|>{\centering\arraybackslash}p{16mm}|>{\centering\arraybackslash}p{16mm}|>{\centering\arraybackslash}p{16mm}|>{\centering\arraybackslash}p{16mm}|}
\hline
\cellcolor{badge!55}\textbf{11} & \cellcolor{badge!55}\textbf{12} & \cellcolor{badge!55}\textbf{13} & \cellcolor{badge!55}\textbf{14}\\
\hline
1 & \textbf{D} & 9 & \textbf{A}\\
\hline
\end{tabular}\\[1.5mm]
\vspace{3.5mm}
{\sffamily\bfseries\color{accent}\fontsize{11}{13}\selectfont Set 3 \ \textmd{\itshape\small Models in context: systems, rates, tables}}\\[1.3mm]
\begin{tabular}{|>{\centering\arraybackslash}p{16mm}|>{\centering\arraybackslash}p{16mm}|>{\centering\arraybackslash}p{16mm}|>{\centering\arraybackslash}p{16mm}|>{\centering\arraybackslash}p{16mm}|>{\centering\arraybackslash}p{16mm}|>{\centering\arraybackslash}p{16mm}|>{\centering\arraybackslash}p{16mm}|>{\centering\arraybackslash}p{16mm}|>{\centering\arraybackslash}p{16mm}|}
\hline
\cellcolor{badge!55}\textbf{1} & \cellcolor{badge!55}\textbf{2} & \cellcolor{badge!55}\textbf{3} & \cellcolor{badge!55}\textbf{4} & \cellcolor{badge!55}\textbf{5} & \cellcolor{badge!55}\textbf{6} & \cellcolor{badge!55}\textbf{7} & \cellcolor{badge!55}\textbf{8} & \cellcolor{badge!55}\textbf{9} & \cellcolor{badge!55}\textbf{10}\\
\hline
\textbf{B} & 3 & \textbf{C} & \textbf{D} & \textbf{C} & \textbf{B} & \textbf{A} & 90 & 8 & \textbf{C}\\
\hline
\end{tabular}\\[1.5mm]
\begin{tabular}{|>{\centering\arraybackslash}p{16mm}|>{\centering\arraybackslash}p{16mm}|>{\centering\arraybackslash}p{16mm}|>{\centering\arraybackslash}p{16mm}|}
\hline
\cellcolor{badge!55}\textbf{11} & \cellcolor{badge!55}\textbf{12} & \cellcolor{badge!55}\textbf{13} & \cellcolor{badge!55}\textbf{14}\\
\hline
6 & \textbf{B} & \textbf{C} & \textbf{D}\\
\hline
\end{tabular}\\[1.5mm]
\vspace{3.5mm}
{\sffamily\bfseries\color{accent}\fontsize{11}{13}\selectfont Set 4 \ \textmd{\itshape\small Inequalities, absolute value and identities}}\\[1.3mm]
\begin{tabular}{|>{\centering\arraybackslash}p{16mm}|>{\centering\arraybackslash}p{16mm}|>{\centering\arraybackslash}p{16mm}|>{\centering\arraybackslash}p{16mm}|>{\centering\arraybackslash}p{16mm}|>{\centering\arraybackslash}p{16mm}|>{\centering\arraybackslash}p{16mm}|>{\centering\arraybackslash}p{16mm}|>{\centering\arraybackslash}p{16mm}|>{\centering\arraybackslash}p{16mm}|}
\hline
\cellcolor{badge!55}\textbf{1} & \cellcolor{badge!55}\textbf{2} & \cellcolor{badge!55}\textbf{3} & \cellcolor{badge!55}\textbf{4} & \cellcolor{badge!55}\textbf{5} & \cellcolor{badge!55}\textbf{6} & \cellcolor{badge!55}\textbf{7} & \cellcolor{badge!55}\textbf{8} & \cellcolor{badge!55}\textbf{9} & \cellcolor{badge!55}\textbf{10}\\
\hline
\textbf{B} & 10 & \textbf{A} & 94 & \textbf{C} & 12 & \textbf{D} & \textbf{D} & \textbf{A} & \textbf{A}\\
\hline
\end{tabular}\\[1.5mm]
\begin{tabular}{|>{\centering\arraybackslash}p{16mm}|>{\centering\arraybackslash}p{16mm}|>{\centering\arraybackslash}p{16mm}|>{\centering\arraybackslash}p{16mm}|}
\hline
\cellcolor{badge!55}\textbf{11} & \cellcolor{badge!55}\textbf{12} & \cellcolor{badge!55}\textbf{13} & \cellcolor{badge!55}\textbf{14}\\
\hline
4 & \textbf{D} & \textbf{B} & \textbf{D}\\
\hline
\end{tabular}\\[1.5mm]
\vspace{3.5mm}
{\sffamily\bfseries\color{accent}\fontsize{11}{13}\selectfont Set 5 \ \textmd{\itshape\small Summit: the hardest mixed problems}}\\[1.3mm]
\begin{tabular}{|>{\centering\arraybackslash}p{16mm}|>{\centering\arraybackslash}p{16mm}|>{\centering\arraybackslash}p{16mm}|>{\centering\arraybackslash}p{16mm}|>{\centering\arraybackslash}p{16mm}|>{\centering\arraybackslash}p{16mm}|>{\centering\arraybackslash}p{16mm}|>{\centering\arraybackslash}p{16mm}|>{\centering\arraybackslash}p{16mm}|>{\centering\arraybackslash}p{16mm}|}
\hline
\cellcolor{badge!55}\textbf{1} & \cellcolor{badge!55}\textbf{2} & \cellcolor{badge!55}\textbf{3} & \cellcolor{badge!55}\textbf{4} & \cellcolor{badge!55}\textbf{5} & \cellcolor{badge!55}\textbf{6} & \cellcolor{badge!55}\textbf{7} & \cellcolor{badge!55}\textbf{8} & \cellcolor{badge!55}\textbf{9} & \cellcolor{badge!55}\textbf{10}\\
\hline
\textbf{B} & −1 & \textbf{C} & 5 & \textbf{C} & 27 & \textbf{B} & \textbf{A} & \textbf{C} & 180\\
\hline
\end{tabular}\\[1.5mm]
\begin{tabular}{|>{\centering\arraybackslash}p{16mm}|>{\centering\arraybackslash}p{16mm}|>{\centering\arraybackslash}p{16mm}|>{\centering\arraybackslash}p{16mm}|>{\centering\arraybackslash}p{16mm}|>{\centering\arraybackslash}p{16mm}|>{\centering\arraybackslash}p{16mm}|>{\centering\arraybackslash}p{16mm}|>{\centering\arraybackslash}p{16mm}|>{\centering\arraybackslash}p{16mm}|}
\hline
\cellcolor{badge!55}\textbf{11} & \cellcolor{badge!55}\textbf{12} & \cellcolor{badge!55}\textbf{13} & \cellcolor{badge!55}\textbf{14} & \cellcolor{badge!55}\textbf{15} & \cellcolor{badge!55}\textbf{16} & \cellcolor{badge!55}\textbf{17} & \cellcolor{badge!55}\textbf{18} & \cellcolor{badge!55}\textbf{19} & \cellcolor{badge!55}\textbf{20}\\
\hline
\textbf{D} & \textbf{B} & \textbf{B} & \textbf{A} & \textbf{C} & 18 & \textbf{B} & \textbf{D} & \textbf{A} & 4\\
\hline
\end{tabular}\\[1.5mm]
\vspace{3.5mm}
}
\end{textblock*}
\clearpage
%%AUTOFLOW-SOLUTIONS
\newgeometry{top=34mm,bottom=24mm,left=16mm,right=16mm}
\pagestyle{solstyle}
\setlength{\columnsep}{14mm}\setlength{\columnseprule}{0.4pt}
{\fontsize{10}{13.6}\selectfont
\begin{multicols}{2}
\solkitem{Set 1}{1}{C ($15$)}{No solution means parallel, distinct lines: $\dfrac k5=\dfrac{12}{4}=3$, so $k=15$. The constants do not follow the ratio ($3\cdot7=21\ne9$), so the lines are distinct. ($\frac53$ inverts the ratio; $3$ is the ratio itself.)}{Ratio test (a : b : c)}
\solitem{2}{Ans. 16}{Infinitely many solutions means the equations are multiples of each other. $21=3\cdot7$, so multiply the first by $3$: $6x+21y=3b$. Then $a=6$ and $3b=30$, $b=10$, so $a+b=16$.}{Ratio test (a : b : c)}
\solitem{3}{A ($-8x+22y=-5$)}{Rearrange: $4x-11y=10$. No solution needs the same slope but a different constant. $-8x+22y=-20$ would be the same line ($-2$ times everything); with $-5$ the lines are parallel. ($8x-22y=20$ is the same line; the other two cross it.)}{Which equation could be the second?}
\solitem{4}{D ($\$1.50$)}{Subtract: $2m=6$, so a muffin is \$3.00. Then $4b=15-9=6$ and $b=\$1.50$. (\$3.00 is the muffin; \$6.00 is $4b$ before dividing by 4; \$4.00 comes from subtracting $5$ instead of $5\cdot3$.)}{Subtract the equations}
\solitem{5}{B}{$x=0$ notebooks and $y=120$ binders: $5\cdot120=600$. The first coordinate is the number of notebooks. (Choice A swaps the coordinates: $3\cdot120=360\ne600$.)}{What does the point mean?}
\solitem{6}{A ($-6$)}{Slope of $\ell$ is $\frac{13-5}{6-2}=2$, so $m$ is $y=2x+b$. Through $(1,-4)$: $-4=2+b$, $b=-6$. (2 is the slope; $-4$ is just the given $y$-coordinate.)}{Parallel: same slope}
\solitem{7}{B ($6$)}{The given slope is $-\frac23$, so $q$ has slope $\frac32$. From $(10,6)$ down to $y=0$ is $6$ units, i.e. $6\div\frac32=4$ units left: $x=6$. (Slopes $\frac23$, $-\frac32$, $-\frac23$ give $1$, $14$, $19$.)}{Perpendicular: negative reciprocal}
\solitem{8}{C ($3$)}{$4x-6=14$ gives $x=5$; $4x-6=-14$ gives $x=-2$. The sum is $3$. (Each single case, or the difference $7$, is a trap.)}{Absolute value: two cases}
\solitem{9}{B ($-4$)}{The left side is $7x-3k$, so the $x$-terms already match. The constants must match too: $-3k=12$, so $k=-4$. (Forgetting to divide by $-3$ gives $-12$.)}{Identity: match the constants}
\solitem{10}{Ans. 11}{$6x-15<4x+9$, so $2x<24$ and $x<12$. The inequality is strict, so $12$ itself fails and the greatest integer is $11$.}{Inequality: strict at the boundary}
\solitem{11}{D $(2,3)$}{$(2,3)$: $3>1$ and $5<6$, both true. $(0,6)$ fails $6<6$; $(4,1)$ fails $1>5$; $(6,3)$ fails both.}{Test each point in both inequalities}
\solitem{12}{Ans. 11/4}{Slope $=\frac{-13-19}{6-(-2)}=-4$, so $f(x)=-4x+11$ (using $f(-2)=8+11=19$). Then $-4x+11=0$ gives $x=\frac{11}4$. (Enter 11/4 or 2.75.)}{Linear function from two points}
\solkitem{Set 2}{1}{C ($6$)}{Parallel needs $\frac k4=\frac9k$, so $k^2=36$ and $k=\pm6$. For $k=-6$ the second equation is $-\frac23$ times the first (the same line: infinitely many solutions). For $k=6$, scaling the first by $\frac23$ gives $4x+6y=4\ne-4$: parallel, no solution.}{Ratio test (a : b : c)}
\solitem{2}{Ans. 5.1}{The equations must be the same line. Multiply the first by $\frac92$ so that $2$ becomes $9$: $\frac32x+\frac{18}5y=9$. Then $p+q=\frac32+\frac{18}5=\frac{51}{10}=5.1$.}{Ratio test (a : b : c)}
\solitem{3}{A ($-10x+8y=-5$)}{Rearrange: $5x-2y-6=2y+14$ gives $5x-4y=20$. No solution needs coefficients in the ratio $5:-4$ but a different constant. $-10x+8y=-5$ is $-2$ times the left side, yet $-2\cdot20=-40\ne-5$: parallel. ($10x-8y=40$ is the same line; $4x+5y=20$ and $5x+4y=20$ cross it.)}{Which equation could be the second?}
\solitem{4}{D ($6x-5y=30$)}{Multiply by the common denominator $15$: $6x-5y=30$. That is the same line, so there are infinitely many solutions. ($5x-6y=30$ swaps the coefficients, $6x-5y=15$ is parallel, $6x+5y=30$ has a different slope.)}{Clear the fractions}
\solitem{5}{B ($8$)}{Put $x=2$ in the second equation: $10+2by=18$, so $by=4$. Then the first gives $2a+4=20$ and $a=8$. (Taking $by=9$ gives 5.5; $by=8$ gives 6; ignoring $by$ gives 10.)}{Substitute the known coordinate}
\solitem{6}{C ($12$)}{Find $\lambda,\mu$ with $\lambda(5,2)+\mu(3,-4)=(1,3)$: $\lambda=\frac12$, $\mu=-\frac12$. So $a+3b=\frac12(34)-\frac12(10)=12$. (Adding the halves gives 22; adding the totals gives 44; $a=6$ alone is not asked.)}{Solve for the combination}
\solitem{7}{A ($4.5$)}{Add: $10x+10y=90$, so $x+y=9$ and the mean is $\frac92=4.5$. (9 is the sum; 90 is the sum of the right sides.)}{Add the equations}
\solitem{8}{B ($2$)}{Parallel needs $\frac k1=\frac4k$, so $k=\pm2$. For $k=2$: $2x+4y=7$ versus $2x+4y=6$ (no solution). For $k=-2$: $-2x+4y=7$ versus $-2x+4y=-6$ (no solution). Both values work, so the answer is 2.}{Ratio test (a : b : c)}
\solitem{9}{Ans. −16}{The two lines that do not contain $k$ meet where $x-y=2$ and $y=2x$: $x-2x=2$, so $x=-2$, $y=-4$. Then $k=2(-2)+3(-4)=-16$.}{Use the extra condition first}
\solitem{10}{C ($5$)}{Multiply the second equation by $-2$ to match $-6y$: $-4x-6y=-2b$. Then $a=-4$ and $-2b=18$, so $b=-9$ and $a-b=5$. ($-13=a+b$, $-5=b-a$, $36=ab$.)}{Ratio test (a : b : c)}
\solitem{11}{Ans. 1}{$3x-2y=5$ is $y=\frac32x-\frac52$. No solution means $\ell$ has the same slope $\frac32$ and a different intercept. Through $(4,7)$: $7=6+b$, so $b=1$.}{Slope-intercept comparison}
\solitem{12}{D ($3x-2y=8$ and $-6x+4y=-16$)}{In $3x-2y=8$ and $-6x+4y=-16$ every number is multiplied by $-2$: the same line. In the first system doubling the second gives $2x+6y=20\ne10$; in the second, doubling the first gives $10x+2y=8\ne6$ (parallel); the last has different slopes.}{Ratio test (a : b : c)}
\solitem{13}{Ans. 9}{Substitute $y$: $2(3x+k)=6x+k+9$, so $6x+2k=6x+k+9$. The $x$-terms cancel; $2k=k+9$ gives $k=9$. (For any other $k$ the lines are parallel and distinct.)}{Substitute and compare}
\solitem{14}{A ($a<b$)}{Set equal: $2x+a=-x+b$, so $x=\frac{b-a}3$. Then $x>0$ means $b>a$. (Try $a=-5,b=-2$ and $a=-1,b=3$ to see that $a+b>0$ and $ab>0$ can fail.)}{Solve in terms of the constants}
\solkitem{Set 3}{1}{B ($30$)}{$p+q=90$ and $0.2p+0.5q=0.3(90)=27$. Substituting $p=90-q$: $18+0.3q=27$, so $q=30$. (60 is the $20\%$ amount; 45 assumes an even split; 27 is the liters of pure acid.)}{Mixture: two equations}
\solitem{2}{Ans. 3}{Upstream speed $=36\div3=12$ and downstream speed $=36\div2=18$. With boat speed $v$ and current $c$: $v-c=12$ and $v+c=18$. Subtract: $2c=6$, so $c=3$.}{Rate equations (add / subtract)}
\solitem{3}{C (It decreases by $3$)}{5 more boxes add $5\cdot18=90$ kg, so the crates must remove $90$ kg: $90\div30=3$ crates. ($\frac35$ is the change per single box; $\frac{25}3$ inverts the ratio.)}{Interpret the coefficients}
\solitem{4}{D ($-9$)}{Slope $=\frac{18-9}{5-2}=3$, so $f(x)=3x+3$. Shifting right 4 gives $g(x)=f(x-4)=3x-9$, with $y$-intercept $-9$. (Shifting left gives 15; shifting up or down by 4 gives 7 or $-1$.)}{Linear function from a table, then shift}
\solitem{5}{C ($\$62$)}{Each extra kilogram costs $\frac{44-26}{7-4}=\$6$, and both parcels are already above 2 kg. Ten kilograms is 3 kg more than seven: $44+3(6)=\$62$. (The flat fee is \$14, so \$74 charges all 10 kg; \$60 drops the fee; \$56 counts only 7 extra kilograms.)}{Slope from two points}
\solitem{6}{B ($80$)}{Add: $50s+50l=400$, so $s+l=8$ and $10s+10l=80$. There is no need to find the prices (\$3 and \$5). (8 is one of each; 400 is the two totals added; 200 is their mean.)}{Add the equations}
\solitem{7}{A}{Turn pipe B off ($y=0$): $\frac x{12}=1$, so $x=12$ minutes. Pipe A alone needs 12 minutes. (Pipe A fills $\frac1{12}$ of the tank per minute, the reciprocal; 20 is pipe B's time.)}{Interpret by setting the other variable to 0}
\solitem{8}{Ans. 90}{Let $u$ be the uphill hours: $18(3.5-u)+12u=54$, so $63-6u=54$ and $u=1.5$ hours. In minutes: $1.5\cdot60=90$.}{Rate equation (watch the units)}
\solitem{9}{Ans. 8}{Slope $=\frac{88-52}{5-2}=12$ per month, so $S=12t+28$. $12t+28>120$ gives $t>7.67$, so the least whole month is $8$. (At $t=7$ the model gives 112, which is still below 120.)}{Linear model, then inequality}
\solitem{10}{C ($40$)}{A linear function changes by the same amount for equal steps: slope $=\frac{12}3=4$. Ten units of input change it by $4\cdot10=40$. (4 is the slope; $12+10=22$ and $12\cdot10=120$ are not proportional reasoning.)}{Slope = change in output per unit input}
\solitem{11}{Ans. 6}{With $c$ correct and $w$ incorrect: $c+w=30$ and $4c-w=90$. Add: $5c=120$, so $c=24$ and $w=30-24=6$.}{Add the equations}
\solitem{12}{B ($28$)}{$M=2J+4$ and $(M+6)+(J+6)=52$. Substitute: $3J+16=52$, so $J=12$ and $M=28$. (12 is Jon's age; 34 is Maya in 6 years; 36 comes from forgetting the 12 extra years.)}{Translate into two equations}
\solitem{13}{C ($125$)}{$28+0.12g=18+0.20g$, so $10=0.08g$ and $g=125$. (50 divides by 0.20 only; 575 adds the fees instead of subtracting; 12.5 slips a decimal place.)}{Set the two costs equal}
\solitem{14}{D ($0$)}{$f(g(x))=2(-x+5)+3=-2x+13$ and $g(f(x))=-(2x+3)+5=-2x+2$. The $x$-terms match but $13\ne2$, so the equation has no solution: 0 values.}{Compose, then compare the equation}
\solkitem{Set 4}{1}{B ($3$)}{Case 1: $x-4=2x-5$ gives $x=1$, but then $|{-3}|=3\ne-3$: extraneous. Case 2: $x-4=-(2x-5)$ gives $x=3$, and $|{-1}|=1=2(3)-5$ works. Only $x=3$, so the sum is 3 (4 would include the fake root).}{Absolute value: check each case}
\solitem{2}{Ans. 10}{$-9\le2x-7\le9$ gives $-2\le2x\le16$, so $-1\le x\le8$. The endpoints are included: $-1,0,\dots,8$ is $10$ integers.}{Absolute value inequality}
\solitem{3}{A $(30,15)$}{$(30,15)$: $45\le50$ and $180+60=240\ge220$. $(10,35)$ earns only $200$; $(25,15)$ earns only $210$; $(35,20)$ uses $55>50$ items.}{Test each point in both inequalities}
\solitem{4}{Ans. 94}{$0.7(74)+0.3f\ge80$ gives $51.8+0.3f\ge80$, so $0.3f\ge28.2$ and $f\ge94$. The least score is 94.}{Weighted average inequality}
\solitem{5}{C ($-1$)}{$(a-5)x<18$. The solution is $x>-3$ (the inequality flipped), so $a-5<0$ and $\frac{18}{a-5}=-3$, giving $a-5=-6$ and $a=-1$. (Ignoring the flip gives 11; adding $5x$ instead of subtracting gives $-11$; combining the constants as $-12+6$ gives 7.)}{Inequality: flip the sign when dividing by a negative}
\solitem{6}{Ans. 12}{The left side is $\frac k3x-3$. Matching $x$-terms: $\frac k3=5$, so $k=15$. Matching constants: $-3=m$. Then $k+m=12$.}{Identity: match coefficients and constants}
\solitem{7}{D (No solution)}{The left side is $5x-10-3x-12=2x-22$. Then $2x-22=2x-20$ becomes $-22=-20$, which is false for every $x$: no solution.}{Contradiction: variables cancel}
\solitem{8}{D ($4.5$)}{The slopes are $-\frac k3$ and $\frac{4}{6}=\frac23$. Perpendicular means the product is $-1$: $-\frac{2k}{9}=-1$, so $k=\frac92$. (Parallel would give $-2$; equal size of slopes gives 2.)}{Perpendicular: negative reciprocal}
\solitem{9}{A ($8$)}{The lines meet where $3x+2=-x+10$: $x=2$, $y=8$. The base on the $y$-axis runs from $2$ to $10$, so it is $8$, and the height is the $x$-coordinate $2$. Area $=\frac12(8)(2)=8$. (16 forgets the $\frac12$; 12 adds the intercepts; 32 uses the $y$-coordinate 8 as the height.)}{Triangle on the y-axis: base along the axis}
\solitem{10}{A ($17$)}{With $b$ large crates, at least $2b$ small crates are needed: $30(2b)+75b=135b\le2400$, so $b\le17.8$ and $b=17$. (18 rounds up; 24 ignores the small crates; 32 is $2400\div75$.)}{Constraints: use the boundary}
\solitem{11}{Ans. 4}{$6<3k-3$ gives $k>3$ (strict), and $k+6\le13$ gives $k\le7$. So $k=4,5,6,7$: four values.}{System of inequalities: intersect the ranges}
\solitem{12}{D ($15$)}{$|x-6|\le4$ means $2\le x\le10$; $|y+2|\le3$ means $-5\le y\le1$. To make $x-y$ largest take $x=10$ and $y=-5$: $15$. (9 uses $y=1$; 11 uses $x=6$; 13 forgets the shift in $y$.)}{Absolute value as an interval}
\solitem{13}{B ($|d-14.5|\le0.3$)}{The midpoint is $\frac{14.2+14.8}2=14.5$ and the half-width is $\frac{14.8-14.2}2=0.3$, so $|d-14.5|\le0.3$. (0.6 is the full width; 14.8 is an endpoint, not the center.)}{Absolute value = center and half-width}
\solitem{14}{D (Infinitely many)}{$|x+2|+|x-4|$ is the sum of the distances from $x$ to $-2$ and to $4$. The points are 6 apart, so every $x$ between them (and including them) gives exactly 6: infinitely many solutions.}{Absolute value: distance reading}
\solkitem{Set 5}{1}{B ($-2$)}{Parallel needs $k(k-1)=2\cdot3$, so $k^2-k-6=0$ and $k=3$ or $-2$. At $k=3$ both equations are $3x+2y=10$ (infinitely many solutions). At $k=-2$ they are $x-y=-5$ and $x-y=\frac{10}3$: parallel and distinct, so no solution.}{Ratio test (two candidates, check the constants)}
\solitem{2}{Ans. −1}{Slope of $\ell$ is $\frac{17-5}{6k-2k}=\frac3k$; the given line has slope $\frac32$, so $k=2$. The points are $(4,5)$ and $(12,17)$: $y=\frac32x+b$ with $5=6+b$, so $b=-1$.}{Slope from two points; parallel}
\solitem{3}{C ($2$)}{$2cx-6=4x-c$. No solution needs equal $x$-coefficients: $2c=4$, $c=2$, and then the constants $-6$ and $-2$ differ. ($c=6$ makes the constants equal but then $x=0$ is a solution.)}{Contradiction: equal x-coefficients, unequal constants}
\solitem{4}{Ans. 5}{Double the first: $4x+2y=6k$; add the second: $5x=7k+5$. Then $y=3k-2x=\frac{k-10}5$ and $x+y=\frac{8k-5}5=7$, so $k=5$. (Check: $x=8$, $y=-1$.)}{Three equations, three unknowns}
\solitem{5}{C ($10$)}{Salt: $0.35(40)=14$ liters stays constant. $14=0.28(40+w)$ gives $40+w=50$ and $w=10$. (14 is the salt amount; 7 is the percent difference; 8 divides by $35\%$.)}{Mixture: the amount of salt does not change}
\solitem{6}{Ans. 27}{For $y=1,2,3,4,5,6$ the inequality $2x\le20-3y$ allows $x\le8,7,5,4,2,1$. That is $8+7+5+4+2+1=27$ pairs. ($y=7$ would need $x\le-0.5$.)}{Lattice points: count row by row}
\solitem{7}{B ($5$)}{$f(f(x))=a(ax+b)+b=a^2x+ab+b$. So $a^2=9$ with $a>0$ gives $a=3$, and $3b+b=8$ gives $b=2$. Then $f(1)=3+2=5$. ($a=-3$ would give $-7$ but is excluded; 11 takes $b=8$.)}{Compose, then match coefficients}
\solitem{8}{A ($4$)}{$|x-3|-2=1$ or $-1$, so $|x-3|=3$ or $|x-3|=1$. That gives $x=0,6$ and $x=2,4$: four solutions. (Stopping after the first layer gives 2.)}{Absolute value: peel the layers}
\solitem{9}{C ($44$)}{Student tickets are cheaper, so buy the most allowed: $y=8$. Then $x\ge4$ to reach 12 tickets: $5(4)+3(8)=44$. (36 uses 12 student tickets, ignoring the cap; 39 buys only 11 tickets; 60 is all adult.)}{Constraints: spend on the cheaper item up to its limit}
\solitem{10}{Ans. 180}{Let $t$ be Car A's hours: $48t=60\left(t-\frac34\right)$, so $12t=45$ and $t=3.75$. Distance $=48(3.75)=180$ miles. (Check: Car B drives 3 hours at 60 mph.)}{Rate equation: equal distances}
\solitem{11}{D}{The lines meet at $m=12$. Plan A starts lower (\$40 vs \$100) but grows faster (\$12 vs \$7 per month), so it is cheaper before $m=12$ and costs more after it. Only the first statement is true.}{Interpret the intersection point}
\solitem{12}{B ($-2$)}{At $x=4$ the $y$-values agree: $4a+3=4b-5$, so $4a-4b=-8$ and $a-b=-2$. (Forgetting to divide by 4 gives $-8$; $b-a$ gives 2; $4(b-a)$ gives 8.)}{Solve for the combination}
\solitem{13}{B ($4.5$)}{Rates per hour: $\frac16+\frac19-\frac1{18}=\frac{3+2-1}{18}=\frac29$ of the tank. Time $=\frac92=4.5$ hours. (3.6 ignores the drain; 7.5 averages 6 and 9.)}{Rates add (work-rate equation)}
\solitem{14}{A ($3x+4y=16$)}{$4x-3y=12$ has slope $\frac43$, so $m$ has slope $-\frac34$. The line $2x+5y=20$ crosses the $y$-axis at $(0,4)$, so $y=-\frac34x+4$, i.e. $3x+4y=16$. ($3x+4y=-16$ uses $-4$, the $y$-intercept of the first line; $3x-4y=-16$ has slope $+\frac34$; $4x+3y=12$ has slope $-\frac43$.)}{Perpendicular + intercept}
\solitem{15}{C ($15$)}{$3x-a=\pm12$ gives $x=\frac{a+12}3$ and $x=\frac{a-12}3$. Their sum is $\frac{2a}3=10$, so $a=15$. (Forgetting to divide by 3 gives 5.)}{Absolute value: two cases}
\solitem{16}{Ans. 18}{Multiply by 3: $-21<5-2x\le15$. Subtract 5: $-26<-2x\le10$. Divide by $-2$ and reverse: $-5\le x<13$. The integers $-5,\dots,12$ number $18$.}{Compound inequality: flip when dividing by a negative}
\solitem{17}{B ($88$)}{Total points: $60(84)=5040$. Class A: $24(78)=1872$. Class B: $5040-1872=3168$, so its mean is $\frac{3168}{36}=88$. (90 averages 78 and 90 equally; 93 swaps the class sizes.)}{Weighted average: use totals}
\solitem{18}{D ($m>-\frac13$)}{$mx+3=-x+9$ gives $x=\frac6{m+1}$ and $y=9-x$. $x>0$ needs $m>-1$; $y>0$ needs $x<9$, so $\frac6{m+1}<9$ and $m>-\frac13$. Both hold exactly when $m>-\frac13$. ($m>-1$ checks only $x>0$; $m>0$ is too strict; $m<-1$ makes $x$ negative.)}{Solve for the intersection, then apply the signs}
\solitem{19}{A ($17$)}{$2(2x+y)=5(x-3y)$ gives $4x+2y=5x-15y$, so $x=17y$ and $\frac xy=17$. (1/17 is $\frac yx$; $-13$ and $-17$ come from sign slips when moving terms.)}{Cross-multiply, then compare coefficients}
\solitem{20}{Ans. 4}{Slope $=\frac{(2p+17)-(2p+1)}{4}=4$, so $y=4x-7$. The first point gives $2p+1=4p-7$, so $p=4$. (Check: $4(8)-7=25=2(4)+17$.)}{Slope from a table with parameters}
\end{multicols}}
\end{document}
